\chapterbackmatter{Supplementary Exercises}

\begin{enumerate}
  \SupplementaryExercise{1}{Vierbein Gauge Counting}
    {(Adapted from Freedman and Van Proeyen, \emph{Ingredients of
    Supergravity}, Sections 2--3)}
    {freedman-van-proeyen-ingredients-sugra-ch31}
  A vierbein has sixteen real components.  Show how local Lorentz
  transformations remove six and identify the remaining ten with a
  symmetric metric.  Linearize the relation
  \(g_{\mu\nu}=e^a{}_\mu e^b{}_\nu\eta_{ab}\).

  \SupplementaryExercise{2}{Spinor Curvature Commutator}
    {(Adapted from Freedman and Van Proeyen, \emph{Ingredients of
    Supergravity}, Sections 2--3)}
    {freedman-van-proeyen-ingredients-sugra-ch31}
  Starting from
  \(D_\mu=\partial_\mu+[\gamma_a,\gamma_b]\omega_\mu{}^{ab}/8\),
  calculate \([D_\mu,D_\nu]\chi\) and express the answer in terms of
  \(R_{\mu\nu}{}^{ab}\).  Check covariance under a local Lorentz
  transformation.

  \SupplementaryExercise{3}{Rarita--Schwinger Gauge Invariance}
    {(Adapted from Freedman and Van Proeyen, \emph{Ingredients of
    Supergravity}, Sections 3--4)}
    {freedman-van-proeyen-ingredients-sugra-ch31}
  For
  \(\mathcal L_{\rm RS}=-\bar\psi_\mu
  i\gamma_5\gamma_\nu\partial_\rho\psi_\sigma
  \epsilon^{\mu\nu\rho\sigma}/2\), prove invariance, up to a
  divergence, under \(\psi_\mu\to\psi_\mu+\partial_\mu\varepsilon\).
  State why the proof changes in a curved background.

  \SupplementaryExercise{4}{Massless Gravitino Degrees of Freedom}
    {(Adapted from Freedman and Van Proeyen, \emph{Ingredients of
    Supergravity}, Sections 3--4)}
    {freedman-van-proeyen-ingredients-sugra-ch31}
  Impose the free equations, the gauge \(\gamma^\mu\psi_\mu=0\), and
  the residual Rarita--Schwinger gauge freedom.  Show that a Majorana
  vector-spinor has two physical helicity states, \(+3/2\) and \(-3/2\).

  \SupplementaryExercise{5}{Super-Higgs Degree Count}
    {(Adapted from Freedman and Van Proeyen, \emph{Ingredients of
    Supergravity}, Sections 3--4)}
    {freedman-van-proeyen-ingredients-sugra-ch31}
  A massless gravitino and a goldstino carry respectively two and two
  on-shell degrees of freedom.  Show that after the goldstino is gauged
  away the resulting massive spin-\(3/2\) field has four.  Identify its
  helicities.

  \SupplementaryExercise{6}{Minkowski Gravitino-Mass Relation}
    {(Adapted from Freedman and Van Proeyen, \emph{Ingredients of
    Supergravity}, Sections 3--4)}
    {freedman-van-proeyen-ingredients-sugra-ch31}
  In a vacuum with only \(F\)-term breaking, use cancellation of the
  vacuum energy to derive the relation between the norm of the order
  parameter \(F\), \(\kappa\), and the gravitino mass.  Reconcile your
  normalization with Eq.~\eqref{eq:31.6.48}.

  \SupplementaryExercise{7}{Torsion from a Gravitino}
    {(Adapted from Van Proeyen, \emph{Tools for Supersymmetry},
    Sections 2.3--2.5)}
    {van-proeyen-tools-susy-ch31}
  Vary a first-order Einstein--Rarita--Schwinger action with respect to
  \(\omega_\mu{}^{ab}\).  Show that the resulting torsion is bilinear
  in \(\psi_\mu\) and explain why substituting its solution generates
  four-gravitino terms.

  \SupplementaryExercise{8}{Why One-and-a-Half Order Works}
    {(Adapted from Van Proeyen, \emph{Tools for Supersymmetry},
    Sections 2.3--2.5)}
    {van-proeyen-tools-susy-ch31}
  Let \(I[e,\psi,\omega(e,\psi)]\) be an action in which \(\omega\)
  solves its own algebraic field equation.  Prove by the chain rule that
  a symmetry check needs no explicit calculation of
  \(\delta\omega(e,\psi)\).

  \SupplementaryExercise{9}{Local-Supersymmetry Closure}
    {(Adapted from Van Proeyen, \emph{Tools for Supersymmetry},
    Sections 2.4--2.5 and 4)}
    {van-proeyen-tools-susy-ch31}
  At the lowest nontrivial order, evaluate
  \([\delta_{\alpha_1},\delta_{\alpha_2}]e^a{}_\mu\) using
  Eqs.~\eqref{eq:31.6.1} and \eqref{eq:31.6.2}.  Identify the general
  coordinate transformation parameter and the terms that must be
  completed by a local Lorentz transformation.

  \SupplementaryExercise{10}{Off-Shell Balance}
    {(Adapted from Van Proeyen, \emph{Tools for Supersymmetry},
    Sections 3--4)}
    {van-proeyen-tools-susy-ch31}
  Count the off-shell bosonic and fermionic components of old-minimal
  \(N=1\) supergravity after gauge redundancies are removed.  Show why
  the auxiliary fields \(s,p,b_\mu\) are required for equality.

  \SupplementaryExercise{11}{AdS Supersymmetry Scale}
    {(Adapted from Van Proeyen, \emph{Tools for Supersymmetry},
    Sections 4.2--4.4)}
    {van-proeyen-tools-susy-ch31}
  Suppose a supersymmetric vacuum has
  \(L_n=0\) but \(f\ne0\).  Use Eq.~\eqref{eq:31.6.68} to find its
  vacuum energy and curvature scale.  Explain why unbroken local
  supersymmetry does not require a flat vacuum.

  \SupplementaryExercise{12}{No-Scale Cancellation}
    {(Adapted from Van Proeyen, \emph{Tools for Supersymmetry},
    Sections 2.4 and 4.3)}
    {van-proeyen-tools-susy-ch31}
  For one neutral field and
  \(\mathcal D=-3\kappa^{-2}\ln(h+h^*)\), verify directly from
  Eq.~\eqref{eq:31.6.77} that the scalar potential vanishes.  State the
  condition needed for a positive kinetic metric.

  \SupplementaryExercise{13}{Anomaly-Mediated Gaugino Mass}
    {(Adapted from Giudice, Luty, Murayama, and Rattazzi,
    \emph{Gaugino Mass without Singlets}, Sections 2--3)}
    {glmr-gaugino-mass-1998-ch31}
  Treat the renormalization scale as a chiral compensator spurion and
  derive
  \(M_\lambda=(\beta_g/g)m_g\), up to the common phase convention.
  Why does this contribution survive without a hidden-sector singlet?

  \SupplementaryExercise{14}{One-Loop Gauge Example}
    {(Adapted from Giudice, Luty, Murayama, and Rattazzi,
    \emph{Gaugino Mass without Singlets}, Sections 2--3)}
    {glmr-gaugino-mass-1998-ch31}
  For \(\beta_g=bg^3/(16\pi^2)\), evaluate the anomaly-mediated
  gaugino mass.  Compare its loop order and parametric size with a
  tree-level mass \(F^i\partial_i f_{AB}\).

  \SupplementaryExercise{15}{Anomaly-Mediated Trilinear Term}
    {(Adapted from Giudice, Luty, Murayama, and Rattazzi,
    \emph{Gaugino Mass without Singlets}, Sections 2--4)}
    {glmr-gaugino-mass-1998-ch31}
  Let \(W=y_{ijk}\Phi_i\Phi_j\Phi_k/6\).  Use wavefunction
  renormalization to show that the anomaly-mediated trilinear
  coefficient is proportional to
  \(-(\gamma_i+\gamma_j+\gamma_k)m_g y_{ijk}\).

  \SupplementaryExercise{16}{Scalar Mass from Running}
    {(Adapted from Giudice, Luty, Murayama, and Rattazzi,
    \emph{Gaugino Mass without Singlets}, Sections 3--4)}
    {glmr-gaugino-mass-1998-ch31}
  Show that the compensator dependence of a wavefunction factor
  \(Z_i(\mu)\) produces
  \(m_i^2=-|m_g|^2\,d\gamma_i/d\ln\mu\,/4\) in the convention
  \(\gamma_i=d\ln Z_i/d\ln\mu\).  Explain the convention dependence of
  the overall sign and factor.

  \SupplementaryExercise{17}{Ultraviolet Insensitivity}
    {(Adapted from Giudice, Luty, Murayama, and Rattazzi,
    \emph{Gaugino Mass without Singlets}, Sections 3--4)}
    {glmr-gaugino-mass-1998-ch31}
  A supersymmetric multiplet of mass \(M\) is integrated out.  Show that
  the threshold change of \(\beta_g\) is accompanied by a matching
  correction to \(M_\lambda\), leaving the low-energy expression
  dependent only on the low-energy beta function.

  \SupplementaryExercise{18}{The Slepton-Sign Problem}
    {(Adapted from Giudice, Luty, Murayama, and Rattazzi,
    \emph{Gaugino Mass without Singlets}, Sections 4--5)}
    {glmr-gaugino-mass-1998-ch31}
  Neglect lepton Yukawa couplings and use the signs of the electroweak
  beta functions to explain why minimal anomaly mediation gives
  negative slepton squared masses.  Why is this a genuine instability
  rather than merely a negative Lagrangian mass parameter?

  \SupplementaryExercise{19}{Kähler-Transformation Invariance}
    {(Adapted from Brignole, Ibanez, and Munoz, \emph{Soft
    Supersymmetry-Breaking Terms from Supergravity and Superstring
    Models}, Section 2)}
    {brignole-ibanez-munoz-soft-terms-ch31}
  Under
  \(d\to d+a+a^*\) and \(f\to fe^{-\kappa^2a}\), prove that
  \(e^{\kappa^2d}|f|^2\) and
  \(e^{\kappa^2d}g^{n\bar m}L_nL_{\bar m}^*\) are invariant.

  \SupplementaryExercise{20}{Minkowski \(F\)-Term Sum Rule}
    {(Adapted from Brignole, Ibanez, and Munoz, \emph{Soft
    Supersymmetry-Breaking Terms from Supergravity and Superstring
    Models}, Sections 2.1--2.2)}
    {brignole-ibanez-munoz-soft-terms-ch31}
  Define
  \(F^n=-e^{\kappa^2d/2}g^{n\bar m}L_{\bar m}^*\).
  Derive the condition on \(g_{n\bar m}F^nF^{*\bar m}\) and \(m_g\)
  in a stationary Minkowski vacuum with vanishing \(D\)-terms.

  \SupplementaryExercise{21}{Tree-Level Gaugino Mass Matrix}
    {(Adapted from Brignole, Ibanez, and Munoz, \emph{Soft
    Supersymmetry-Breaking Terms from Supergravity and Superstring
    Models}, Section 2.2)}
    {brignole-ibanez-munoz-soft-terms-ch31}
  From Eq.~\eqref{eq:31.6.75}, derive the canonically normalized
  tree-level gaugino mass for diagonal gauge kinetic function
  \(f_A(\phi)\).  Show explicitly why a constant \(f_A\) gives zero.

  \SupplementaryExercise{22}{Visible Scalar Soft Mass}
    {(Adapted from Brignole, Ibanez, and Munoz, \emph{Soft
    Supersymmetry-Breaking Terms from Supergravity and Superstring
    Models}, Sections 2.2--2.3)}
    {brignole-ibanez-munoz-soft-terms-ch31}
  A visible field \(Q\) has Kähler metric \(Z_Q(z,z^*)\) and vanishing
  vacuum expectation value, while hidden fields \(z^m\) break
  supersymmetry.  Derive
  \(m_Q^2=m_g^2-F^mF^{*\bar n}
  \partial_m\partial_{\bar n}\ln Z_Q\) for zero vacuum energy.

  \SupplementaryExercise{23}{Gravity-Mediated \(A\)-Term}
    {(Adapted from Brignole, Ibanez, and Munoz, \emph{Soft
    Supersymmetry-Breaking Terms from Supergravity and Superstring
    Models}, Sections 2.2--2.3)}
    {brignole-ibanez-munoz-soft-terms-ch31}
  For a visible Yukawa coupling \(Y_{ijk}(z)\) and diagonal visible
  metrics \(Z_i\), derive the hidden-\(F\) contribution proportional to
  \[
    F^m\partial_m
    \ln\!\left(\frac{e^{\kappa^2d}Y_{ijk}}
      {Z_iZ_jZ_k}\right).
  \]

  \SupplementaryExercise{24}{Flavor Universality Criterion}
    {(Adapted from Brignole, Ibanez, and Munoz, \emph{Soft
    Supersymmetry-Breaking Terms from Supergravity and Superstring
    Models}, Sections 2.3 and 3)}
    {brignole-ibanez-munoz-soft-terms-ch31}
  Use the result of the preceding two exercises to give sufficient
  conditions for generation-universal scalar masses and aligned
  trilinear terms.  Explain which derivatives of the Kähler metric
  threaten flavor conservation.

  \SupplementaryExercise{25}{Planck-Suppressed Mediation Scale}
    {(Adapted from Martin, \emph{A Supersymmetry Primer}, version 7,
    Sections 7 and 12)}
    {martin-susy-primer-v7-ch31}
  A hidden sector has \(\sqrt F\ll\kappa^{-1}\).  By dimensional
  analysis determine the scale of visible soft masses generated by
  operators suppressed by one power of the reduced Planck mass.  Relate
  it to the gravitino mass when the cosmological constant is tuned.

  \SupplementaryExercise{26}{Goldstino Direction}
    {(Adapted from Martin, \emph{A Supersymmetry Primer}, version 7,
    Sections 7.1 and 12.1)}
    {martin-susy-primer-v7-ch31}
  For chiral fermions \(\psi_i\) and gauginos \(\lambda_A\) with vacuum
  auxiliary values \(F_i,D_A\), find the normalized goldstino linear
  combination.  Show that it shifts by a constant under the broken
  supercharge.

  \SupplementaryExercise{27}{Longitudinal-Gravitino Equivalence}
    {(Adapted from Martin, \emph{A Supersymmetry Primer}, version 7,
    Sections 7.1 and 12)}
    {martin-susy-primer-v7-ch31}
  For energy \(E\gg m_g\), write the leading longitudinal
  spin-\(3/2\) polarization in terms of \(q_\mu/m_g\) and a spin-\(1/2\)
  wavefunction.  Explain why high-energy amplitudes with longitudinal
  gravitinos equal goldstino amplitudes up to \(O(m_g/E)\).

  \SupplementaryExercise{28}{Gravitino Decay Scaling}
    {(Adapted from Martin, \emph{A Supersymmetry Primer}, version 7,
    Sections 12 and 13)}
    {martin-susy-primer-v7-ch31}
  Use the goldstino coupling to estimate the two-body decay width of a
  superpartner of mass \(m\) to its partner plus a light gravitino.
  Determine the dependence on \(m\), \(F\), and \(m_g\).

  \SupplementaryExercise{29}{Polonyi Stability Check}
    {(Adapted from Martin, \emph{A Supersymmetry Primer}, version 7,
    Sections 7.1 and 12.1)}
    {martin-susy-primer-v7-ch31}
  For the model of Weinberg Exercise 6, compute the two real scalar
  squared masses at the Minkowski stationary point.  Verify that both
  are positive and express them in units of the gravitino mass.

  \SupplementaryExercise{30}{Hidden-Sector Scale Comparison}
    {(Adapted from Martin, \emph{A Supersymmetry Primer}, version 7,
    Sections 12.1--12.4)}
    {martin-susy-primer-v7-ch31}
  Compare a hidden sector with \(F=O(\Lambda^2)\) to one whose
  nonperturbative superpotential is \(O(\Lambda^3)\).  Derive the
  characteristic gravity-mediated scales
  \(\kappa\Lambda^2\) and \(\kappa^2\Lambda^3\), and solve for
  \(\Lambda\) when the visible scale is \(1\) TeV.
\end{enumerate}
